\subsection{Alien measures}

Given two real measures \(\rho,\sigma\) on T, recall that \(\rho\) and \(\sigma\) are said to be alien (to each other) if \(\inf(|\rho|,|\sigma|)=0\) in \(\mathcal{M}(\mathrm{T})\) (Ch. II, §1, No.~1). The real measures alien to a given measure are known to form a band (Ch. II, §1, No.~5, Th. 1). This definition may be extended immediately to the case of complex measures.

DEFINITION 4. --- One says that a complex measure \(\theta\) on T is concentrated on a subset M of T, or that M carries \(\theta\) , if \(\mathbb{C}M\) is locally negligible for \(\theta\) .

The set M carries \(\theta\) if and only if it carries \(|\theta|\) . It is equivalent to say that M carries \(\theta\) , or that M is \(\theta\) -measurable and \(\theta = \varphi_{M} \cdot \theta\) . If \(\theta\) is concentrated on M, then every measure with base \(|\theta|\) is concentrated on M.

PROPOSITION 13. --- In order that two complex measures \(\rho\) and \(\sigma\) on T be alien to each other, it is necessary and sufficient that there exist in T two disjoint sets R and S such that \(\rho\) is concentrated on R and \(\sigma\) on S; R and S may be taken to be universally measurable.

Set \(\mu = |\rho|\) , \(\nu = |\sigma|\) , \(\lambda = \mu + \nu\) ; since \(\mu\) and \(\nu\) are bounded above by \(\lambda\) , there exist two locally \(\lambda\) -integrable functions u and v (which one can suppose to be universally measurable by §3, No.~4, Prop. 7) such that \(\mu = u \cdot \lambda\) , \(\nu = v \cdot \lambda\) . Then

\[
\inf (| \rho |, | \sigma |) = \inf (\mu , \nu) = \inf (u, v) \cdot \lambda
\]

(No.~2, Cor. of Prop. 2). Let A (resp. B) be the set of \(t \in \mathbf{T}\) such that \(u(t) > 0\) and \(v(t) = 0\) (resp. \(u(t) = 0\) and \(v(t) > 0\) ). If \(\rho\) and \(\sigma\) are alien, then \(\inf(u, v) = 0\) locally \(\lambda\) -almost everywhere (No.~3, Cor. 2 of Prop. 3), so that the disjoint universally measurable sets A and B carry \(\mu\) and \(\nu\) , respectively. Conversely, suppose that \(\mu\) and \(\nu\) are carried, respectively, by disjoint sets R and S; \(\varphi_{\mathrm{R}}\) is measurable for the measure \(\mu = u \cdot \lambda\) , and \(\mu = \varphi_{\mathrm{R}} \cdot \mu\) . By Prop. 8 of No.~4, the function \(u' = u\varphi_{\mathrm{R}}\) is \(\lambda\) -measurable, and \(\mu = u' \cdot \lambda\) . Similarly, if \(v' = v\varphi_{\mathrm{S}}\) then \(\nu = v' \cdot \lambda\) ; one concludes by remarking that \(\inf(u', v') = 0\) (No.~2, Cor. of Prop. 2).

COROLLARY 1. --- For every real measure \(\nu\) on T, there exist two disjoint sets M,N carrying \(\nu^{+}\) and \(\nu^{-}\) , respectively.

\pandocbounded{\includegraphics[keepaspectratio]{images/part002_61ad429f.jpg}}

One must take care not to confuse the notion of support of a measure \(\nu\) , and that of a set where \(\nu\) is concentrated. The support S of \(\nu\) is the smallest closed set carrying \(\nu\) (Ch. III, §2, No.~2, Prop. 2 and Ch. IV, §2, No.~2, Prop. 5). However, there may exist subsets of S that are distinct from S and carry \(\nu\) . More generally, one can have \(\inf(\mu,\nu)=0\) for two positive measures \(\mu\) and \(\nu\) having the same support (Exer. 5).

Note also that the intersection of the sets carrying \(\nu\) is the set of points \(t \in T\) such that \(|\nu|(\{t\}) > 0\) , and it can be empty (for example, in the case of Lebesgue measure); therefore, there is not in general a smallest set carrying \(\nu\) .

COROLLARY 2. --- Let \(\rho\) and \(\sigma\) be two alien complex measures, and let \(\rho'\) and \(\sigma'\) be two complex measures admitting densities relative to \(\rho\) and \(\sigma\) , respectively; then \(\rho'\) and \(\sigma'\) are alien.

COROLLARY 3. --- Let \(\rho\) and \(\sigma\) be two alien complex measures; then \(|\rho + \sigma| = |\rho| + |\sigma|\) .

Denote by v (resp. w) a universally measurable function of absolute value 1 such that \(\rho = v \cdot |\rho|\) (resp. \(\sigma = w \cdot |\sigma|\) ) (Cor. 3 of Th. 2), and by A a universally measurable set carrying \(\rho\) , such that B = C A carries \(\sigma\) (Prop. 13); then also \(\rho + \sigma = (v\varphi_{\mathrm{A}} + w\varphi_{\mathrm{B}}) \cdot (|\rho| + |\sigma|)\) . Since the function \(v\varphi_{A} + w\varphi_{B}\) has absolute value equal to 1, the corollary follows from Prop. 2.

THEOREM 3 (Lebesgue). --- Every complex measure \(\theta\) on T may be written in one and only one way in the form \(\theta = g \cdot \mu + \theta'\) , where \(g\) is locally \(\mu\) -integrable and \(\theta'\) is a measure alien to \(\mu\) . Then \(|\theta| = |g| \cdot \mu + |\theta'|\) .

When \(\theta\) is positive, this follows at once from the theorem of F. Riesz (Ch. II, §1, No.~5, Th. 1) applied to the fully lattice-ordered space \(\mathcal{M}(T)\) of real measures on T, and to the band generated by \(\mu\) in this space, on taking into account Cor. 2 of No.~5, Th. 2; moreover, \(\theta'\) and \(g \cdot \mu\) are then positive, which implies that \(g\) is positive locally \(\mu\) -almost everywhere (Cor. 3 of Prop. 3). To treat the case that \(\theta\) is not positive, set \(\nu = |\theta|\) , \(\nu = f \cdot \mu + \nu'\) (where \(f\) is positive and where \(\nu'\) and \(\mu\) are alien to each other), and \(\theta = v \cdot \nu\) , where \(v\) is a universally measurable function of absolute value 1 (Cor. 3 of Th. 2). We then have (Prop. 8) \(\theta = g \cdot \mu + \theta'\) , with \(g = vf\) (so that \(|g| = f\) ) and \(\theta' = v \cdot \nu'\) (so that \(|\theta'| = \nu'\) by Prop. 2); the measures \(\theta'\) and \(\mu\) are alien to each other by Cor. 2 of Prop. 13. It only remains to establish the uniqueness of the decomposition. Thus, suppose that \(\theta = g \cdot \mu + \theta' = g_1 \cdot \mu + \theta'_1\) , where \(\theta'\) and \(\theta'_1\) are alien to \(\mu\) ; \(|\theta' - \theta'_1|\) is bounded above by \(|\theta'| + |\theta'_1|\) , therefore \(\theta' - \theta'_1\) is alien to \(\mu\) , hence also to \((g_1 - g) \cdot \mu\) . The relation \(\theta' - \theta'_1 = (g_1 - g) \cdot \mu\) then implies that the two members are zero, which proves uniqueness.

Recall (No.~5, Th. 2 and Scholium) that the space \(\mathrm{L}_{\mathrm{loc}}^{1}(\mathrm{T},\mu;\mathbf{C})\) may be identified (by means of the mapping \(g\mapsto g\cdot\mu\) ) with a subspace of \(\mathcal{M}_{\mathbf{C}}(\mathrm{T})\) . With this convention, Theorem 3 takes the following form:

COROLLARY. --- There exists a projector \(p\) of the space \(\mathcal{M}_{\mathbf{C}}(\mathrm{T})\) onto the space \(\mathrm{L}_{\mathrm{loc}}^{1}(\mathrm{T},\mu ;\mathbf{C})\) , whose kernel \(\overline{p}^{1}(0)\) is the set of complex measures alien to \(\mu\) , such that

\[
| \theta | = | p (\theta) | + | \theta - p (\theta) |, \quad p (| \theta |) = | p (\theta) |
\]

for every complex measure \(\theta\) .

If p is restricted to the set of bounded measures, one obtains a projector \(p^{1}\) of the space \(\mathcal{M}_{\mathbf{C}}^{1}(\mathrm{T})\) onto the space \(\mathrm{L}_{\mathbf{C}}^{1}(\mathrm{T},\mu)\) ; the relation \(\|\theta\|=|\theta|(1)\) implies that \(\|\theta\|=\|p^{1}(\theta)\|+\|\theta-p^{1}(\theta)\|\) for every bounded complex measure \(\theta\) .

